\documentclass{../note}
\usepackage{bbold}
\begin{document}

\begin{zadanie}{2}
    Niech $\alpha,\beta\in\mathbb{k}^{n}$ będą wektorami o współrzędnych $\alpha=(a_{1},\dots,a_{n})$, $\beta=(b_{1},\dots,b_{n})$. Oznaczmy przez $M(\alpha,\beta)=(m_{ij})_{i,j\in 1\dots n}$ macierz $n\times n$, której wyrazy są równe $m_{ij}=a_{i}\cdot b_{j}$ dla $i,j\in 1\dots n$.
\begin{enumerate}
    \item[(a)] Udowodnij, że dla dowolnych $\alpha,\beta\in\mathbb{k}^{n}$ rząd macierzy $M(\alpha,\beta)$ wynosi co najwyżej 1.
    \item[(b)] Udowodnij, że każda macierz ze zbioru $M_{n\times n}(\mathbb{k})$ o rzędzie 1 jest równa $M(\alpha,\beta)$ dla pewnych $\alpha,\beta\in\mathbb{k}^n$.
    \item[(c)] Załóżmy że wektory $\alpha,\alpha^{\prime}\in\mathbb{k}^{n}$ są liniowo niezależne oraz, że wektory $\beta, \beta^{\prime}\in\mathbb{k}^{n}$ są liniowo niezależne. Znajdź rząd macierzy $M(\alpha,\beta)+M(\alpha^{\prime},\beta^{\prime})$.
\end{enumerate}
\end{zadanie}
\begin{rozwiazanie}
\begin{enumerate}[label=(\alph*)]
    \item $i$-ty rząd $M(\alpha, \beta)$ to $\beta_i\alpha$. Każdy rząd jest wielokrotnością $\alpha$, czyli każda para rzędów jest liniowo zależna. Zatem wymiar przetrzeni liniowej rozpiętej przez rzędy to maksymalnie 1, czyli rząd macierzy to maksymalnie 1.
    \item Jeśli rząd macierzy to 1, to każda para rzędów jest liniowo zależna, czyli m.in. każdy rząd jest liniowo zależny od pewnego niezerowego rzędu. Niech ten rząd to $\alpha$. Wtedy $i$-ty rząd musi być postaci $\beta_i\alpha$, czyli cała macierz musi być postaci $M(\alpha, \beta)$.
    \item Wiadomo, że $r(M(\alpha, \beta) + M(\alpha', \beta')) \leq r(M(\alpha, \beta)) + r(M(\alpha', \beta')) \leq 1 + 1 = 2$. Załóżmy niewprost, że $r(M(\alpha, \beta) + M(\alpha', \beta')) \leq 1$. Wtedy każda para rzędów jest liniowo zależna, czyli dla każdych $i, j \in [1, n]$ istnieją takie $x, y \in \mathbb{k}$, że zachodzi
    \[x(\beta_i\alpha + \beta'_i\alpha') + y(\beta_j\alpha + \beta'_j\alpha') = 0\]
    Można to przekształcić:
    \[(x\beta_i + y\beta_j)\alpha + (x\beta'_i + y\beta'_j)\alpha' = 0\]
    Skoro $\alpha$ i $\alpha'$ są liniowo niezależne, to 
    \[x\beta_i + y\beta_j = x\beta'_i + y\beta'_j = 0\]
    $\beta \neq 0$ bo jest liniowo niezależne od $\beta'$, czyli istnieje takie $j$, że $\beta_j \neq 0$. Wtedy dla każdego $i$
    \[\frac{\beta_i}{\beta_j} = \frac{-y}{x} = \frac{\beta'_i}{\beta'_j}\]
    Niech $\gamma_i = \frac{\beta_i}{\beta_j}$ oraz $\gamma = (\gamma_1, \ldots \gamma_n)$. Z powyższego równania wiemy, że $\beta = (\beta_j\gamma_1, \ldots, \beta_j\gamma_n) = \beta_j\gamma$ oraz $\beta' = (\beta'_j\gamma_1, \ldots, \beta'_j\gamma_n) = \beta'_j\gamma$. Zatem $\beta$ i $\beta'$ są wielokrotnością tego samego wektora, czyli są liniowo zależne, co jest sprzecznością. Zatem ranga naszej macierzy to 2.
\end{enumerate}
\end{rozwiazanie}

\begin{zadanie}{3}
Udowodnij, że dla dowolnych liczb $a, b, c, d\in\mathbb{C}$ spełniających warunek $ad -bc \ne 0$ zbiór
\[ \{z\in\mathbb{C} : \left|\frac{az+b}{cz+d}\right|=1\} \]
jest okręgiem lub prostą na płaszczyźnie zespolonej. Wykaż, że każdy okrąg i każdą prostą na płaszczyźnie zespolonej można przedstawić jako powyższy zbiór dobierając odpowiednio $a, b, c, d\in\mathbb{C}$
\end{zadanie}
\begin{rozwiazanie}
Zbiór $z \in \C$ spełniających $|az + b| = 1$, czyli $\left|z - \left(-\frac{b}{a}\right)\right| = \frac1a$, to zbiór $z$ odległych od $(-\frac{b}{a})$ o $\left|\frac1a\right|$ na płaszczyźnie zespolonej, czyli okrąg o środku $(-\frac{b}{a})$ i promieniu $\left|\frac1a\right|$. Zatem zbiór $z$ spełniających $\left|\frac{az+b}{cz+d}\right|=1$ to zbiór $z$ o stałym stosunku odległości od pewnych dwóch punktów, a konkretnie to te punkty to $-\frac{b}{a}$ i $-\frac{c}{d}$, a stosunek odległości to $\left|\frac{c}a\right|$. Zamiast pokazywać, że to prawda dla dowolnych punktów $-\frac{b}{a}$ i $-\frac{c}{d}$, możemy powołać się na podobieństwo (w znaczeniu geometrzycznym) i pokazać, że dla $A = (0, 0)$ i $B = (1, 0)$, zbiór punktów $C$ dla których $\frac{BC}{AC} = r$, to okrąg albo prosta. Pokażemy to dla $r \geq 1$, a przeciwny przypadek jest analogiczny.\\
Niech $C = (x, y)$. Gdy $r=1$, to zbiór takich $C$ to prosta. Załóżmy teraz, że $r > 1$, czyli że $r^2 - 1 > 0$. Warunek o stosunku odległości możemy podnieść do kwadratu:
\[\frac{(x-1)^2+y^2}{x^2+y^2} = r^2\]
\[x^2(r^2 - 1) + 2x + \frac{1}{r^2 - 1} + y^2(r^2-1) = 1 + \frac{1}{r^2 - 1}\]
\[\left(x + \frac{1}{r^2 - 1}\right)^2 + y^2 = \frac{r^2}{(r^2 - 1)^2}\]
Zatem wszystkie takie $C$ spełniają równianie pewnego okręgu. Warto też powiedzieć, że to przekształcenia można cofnąć, czyli każde $C$ na tym okręgu spełnia warunek.\\
Pozostało wykazać, że każdy okrąg i każdą prostą można przedstawić w ten sposób. Jeśli mamy okrąg o środku $s$ i promieniu $r$, to możemy wybrać $a$ i $b$ takie, że $r= \frac1a$ i $s = -\frac{b}{a}$, czyli $a = \frac1r$ i $b = -\frac{s}{r}$ oraz $c = 0$ i $d = 1$, żeby mianownik był równy 1. Jeśli mamy prostą, to możemy wybrać takie dwa punkty $-b$ i $-d$, dla których ta prosta jest symetralną oraz $a=c=1$, to wtedy $\left|\frac{az+b}{cz+d}\right|=1$ dla punktów na tej prostej.
\end{rozwiazanie}

\begin{zadanie}{4}
Załóżmy, że $\mathbb{k}$ jest ciałem charakterystyki $p>0$.
\begin{enumerate}
    \item[(a)] Dowiedź, że $(a+b)^{p^{n}}=a^{p^{n}}+b^{p^{n}}$ dla dowolnych $a, b\in\mathbb{k}$ oraz $n\ge 1$.
    \item[(b)] Wywnioskuj z (a), że odwzorowanie $F:\mathbb{k}\rightarrow\mathbb{k}$ dane wzorem $F(x)=x^{p}$ jest morfizmem ciał, to znaczy zachowuje działania, czyli spełnia $F(a+b)=F(a)+F(b)$ oraz $F(a\cdot b)=F(a)\cdot F(b)$ dla dowolnych $a, b\in\mathbb{k}$. Ponadto $F(0)=0$ i $F(1)=1$.
    \item[(c)] Uzasadnij, że odwzorowanie $F$ z punktu (b) jest injektywne. W szczególności gdy ciało $\mathbb{k}$ jest skończone, to $F$ jest bijekcją oraz $\mathbb{k}=\mathbb{k}^{p}=\{a^{p}:a\in\mathbb{k}\}$.
    \item[(d)] Wskaż przykład ciała $\mathbb{k}$, dla którego $\mathbb{k}\ne\mathbb{k}^{p}$, to znaczy takiego, że odwzorowanie $F$ z punktu (b) nie jest surjektywne.
\end{enumerate}
\end{zadanie}
\begin{rozwiazanie}
\begin{enumerate}[label=(\alph*)]
    \item Wiadomo, że ciała skończone są przestrzenią liniową nad $\Z_p$. Niech $m$ to wymiar tej przestrzeni (dowód działa tylko dla skończonych $m$) i niech $(a_k)_{1 \leq k \leq m}$ to jej baza. Najpierw pokażemy, że jeśli $x = \sum_{k=1}^{m} w_ka_k$, to $x^p = \sum_{k=1}^{m} w_ka_k^p$. Wiemy, że to jest prawda dla $x = 0$. Przeprowadźmy dowód indukcyjny po kolejnych wymiarach $k$ i dla każdego wymiaru przeprowadźmy dowód indukcyny po $w_k$. Załóżmy, że dla danego $x$, $k$ to ostatni wymiar, na którym $w_k > 0$. Wtedy
    \[x^p = \left(\sum_{i = 1}^{k} w_ia_i\right)^p = \left(a_k + \sum_{i = 1}^{k} (w_i-\delta_{ik})a_i\right)^p = \sum_{l=0}^{p} \binom{p}{l}a_k^l\left(\sum_{i = 1}^{k} (w_i-\delta_{ik})a_i\right)^{p-l}\]
    Dla $0 < l < p$, $\binom{p}{l} \equiv 0 (\mod p)$, czyli
    \[x^p = a_k^p + \left(\sum_{i = 1}^{k} (w_i-\delta_{ik})a_i\right)^p\]
    Możemy teraz skorzystać z założenia indukcyjnego dla tego wymiaru - wiemy, że 
    \[\left(\sum_{i = 1}^{k} (w_i-\delta_{ik})a_i\right)^p = \sum_{i = 1}^{k} (w_i-\delta_{ik})a_i^p\]
    Zatem
    \[x^p = a_k^p + \sum_{i = 1}^{k} (w_i-\delta_{ik})a_i^p = \sum_{i = 1}^{k} w_ia_i^p\]
    Niech $x = \sum_{k=1}^{m} w_ka_k$ i $y = \sum_{k=1}^{m} v_ka_k$. Wtedy 
    \[(x + y)^p = \left(\sum_{k=1}^{m} (w_k+v_k)a_k\right)^p = \sum_{k=1}^{m} (w_k+v_k)a_k^p = x^p + y^p\]
    Możemy teraz poprowadzić dowód indukcyjny po $n$:
    \[(x + y)^{p^n} = \left((x + y)^p\right)^{p^{n-1}} = \left((x^p + y^p)\right)^{p^{n-1}} = x^{p^n} + y^{p^n}\]
    \item Już wiemy, że $F(x + y) = (x + y)^p = x^p + y^p = F(x) + F(y)$. Wiadomo, też, że $F(xy)^p = (xy)^p = x^py^p = F(x)F(y)$. Zachodzi też $F(0) = 0^p = 0$ oraz $F(1) = 1^p = 1$.
\end{enumerate}
\end{rozwiazanie}

\begin{zadanie}{5}
Niech $\mathbb{R}[x]_{\le 4}\subset\mathbb{R}[x]$ oznacza przestrzeń wielomianów stopnia $\le 4,$ czyli wielomianów postaci $w(x)=a_{0}+a_{1}x+a_{2}x^{2}+a_{3}x^{3}+a_{4}x^{4}$ gdzie $a_{0},\dots,a_{4}\in\mathbb{R}$. Niech $W\subset\mathbb{R}[x]_{\le 4}$ będzie zbiorem tych wielomianów $w$ stopnia 4, dla których funkcje wielomianowe $\tilde{w}:\mathbb{R}\rightarrow\mathbb{R}$ spełniają warunek
\[ \forall_{x\in\mathbb{R}\backslash\{0\}} \quad x^{4}\cdot\tilde{w}\left(\frac{1}{x}\right)=\tilde{w}(x) \]
\begin{enumerate}
    \item[(a)] Przedstaw podzbiór $W\subseteq\mathbb{R}[x]_{\le 4}$ jako zbiór rozwiązań układu równań jednorodnych zadających relacje między współczynnikami $a_{0},\dots,a_{4}\in\mathbb{R}$. Uzasadnij, że $W$ jest przestrzenią liniową nad ciałem $\mathbb{R}$.
    \item[(b)] Rozpatrzmy wielomiany
    \[ w_{1}=1+tx^{2}+x^{4}, \quad w_{2}=sx^{2}, \quad w_{3}=1+tx+x^{2}+tx^{3}+x^{4} \]
    zależne od parametrów $s, t\in\mathbb{R}.$ Jakie warunki muszą spełniać parametry $s$ i $t$, aby wielomiany $w_{1}, w_{2}, w_{3}$ stanowiły bazę przestrzeni $W$?
    \item[(c)] Połóżmy $s=t=1$. Znajdź współrzędne wielomianu $x+2x^{2}+x^{3}$ w bazie składającej się z $w_{1}, w_{2}$ i $w_{3}$.
\end{enumerate}
\end{zadanie}
\begin{rozwiazanie}
\begin{enumerate}[label=(\alph*)]
    \item Jeśli $w(x)=a_{0}+a_{1}x+a_{2}x^{2}+a_{3}x^{3}+a_{4}x^{4}$, to dla $x \neq 0$, to $x^4 \tilde{w}\left(\frac1x\right) = a_{0}x^4+a_{1}x^3+a_{2}x^{2}+a_{3}x^{1}+a_{4}$. Jeśli dwa wielomiany stopnia 4 są równe na $\R \backslash \{0\}$, to są równe m.in. w pewnych 5 punktach, czyli z wzoru interpolacji Lagrange'a muszą mieć równe współczynniki. Zatem każde $w \in W$ ma współczynniki spełniające $a_0 = a_4$ i $a_1 = a_3$. Łatwo sprawdzić, że to jest warunek wystarczający. Jest to więc układ równań określający $W$, którego szukamy. Niech $w(x)=a_{0}+a_{1}x+a_{2}x^{2}+a_{3}x^{3}+a_{4}x^{4}$ i $v(x)=b_{0}+b_{1}x+b_{2}x^{2}+b_{3}x^{3}+b_{4}x^{4}$ to dwa wielomiany należące do $W$. Wtedy jeśli $u(x) = w(x) + v(x) =c_{0}+c_{1}x+c_{2}x^{2}+c_{3}x^{3}+c_{4}x^{4}$, to $c_0 = a_0 + b_0 = a_4 + b_4 = c_4$ oraz $c_1 = a_1 + b_1 = a_3 + b_3 = c_3$, czyli $u \in W$. Jeśli $u$ spełnia $u(x) = dw(x)$, to $c_0 = dc_0 = dc_4 = c_4$ i $c_1 = dc_1 = dc_3 = c_3$. Elementem zerowym jest wielomian zerowy. Mamy więc zdefiniowane operacje definiujące przestrzeń liniową. Łatwo sprawdzić, że wszystkie aksjomaty przestrzeni liniowej są spełnione.
    \item Zauważmy, że $b_1(x) = 1 + x^4, b_2(x) = x + x^3, b_3(x) = x^2$ są bazą $W$, bo dowolne $w \in W$ musi być postaci $ax^4 + bx^3 + cx^2 + bx + a = ab_1 + bb_2 + cb_3$. Widać, że są liniowo niezależne. Wynika z tego, że $\dim(W) = 3$, więc $w_1, w_2, w_3$ muszą być liniowo niezależne. $w_1, w_2, w_3$ mają współrzędne $(1, 0, t), (0, 0, s), (1, t, 1)$ w bazie $(b_1, b_2, b_3)$. Żeby się dowiedzieć, czy są liniowo niezależne, możemy zeschodkować macierz:
    \[
    \begin{pmatrix}
        1 & 0 & t \\
        1 & t & 1 \\
        0 & 0 & s
    \end{pmatrix} \xrightarrow{\text{II - I}}
    \begin{pmatrix}
        1 & 0 & t \\
        0 & t & 1 - t \\
        0 & 0 & s
    \end{pmatrix}
    \]
    Wektory są liniowo niezależne, gdy są wszystkie schodki, czyli gdy $t \neq 0$ i $s \neq 0$.
    \item Łatwo sprawdzić, że wektor o współrzędnych $(-1, 2, 1)$ to właśnie wielomian $x+2x^{2}+x^{3}$.
\end{enumerate}
\end{rozwiazanie}

\end{document}